This independent annex contains supplementary material for the report. The first section provides complete SPARQL queries moved from Chapter~3; additional annex sections may be added here without changing the document's structure. The prefixes declared in the first Chapter~3 listing are assumed unless a listing explicitly introduces a new prefix.

\section{Exploratory Queries}
\label{sec:annex-a1}

\begin{lstlisting}[
	caption={Reverse dependency chain: elements depending on a given component},
	label={lst:annex-reverse-dep-query}
]
SELECT DISTINCT ?component ?service ?serviceName
WHERE {
  ?component rdfs:label "PaaS Services"@en .
  ?component (archimate:DependencyRelationship
             |archimate:StructuralRelationship)+ ?service .
  ?service rdf:type archimate:Element ;
           rdfs:label ?serviceName .
}
\end{lstlisting}

\begin{lstlisting}[
	caption={List elements with their ArchiMate layer},
	label={lst:annex-layer-element-query}
]
# Newly introduced prefix for this listing.
PREFIX owl: <http://www.w3.org/2002/07/owl#>

SELECT ?s ?n ?layer
WHERE {
  ?s a archimate:Element, ?layer .
  ?layer a archimate:Layer .
  ?s rdfs:label ?n .
}
\end{lstlisting}

\begin{lstlisting}[
	caption={Cross-layer relationship counts: Application to Business},
	label={lst:annex-inter-layer-connections-query}
]
SELECT ?src ?target (COUNT(?rel) AS ?counter)
WHERE {
  ?src rdfs:subClassOf archimate:ApplicationLayer .
  ?target rdfs:subClassOf archimate:BusinessLayer .
  ?a a ?src .
  ?b a ?target .
  ?a ?rel ?b .
}
GROUP BY ?src ?target
\end{lstlisting}

\section{Gap Analysis Queries}
\label{sec:annex-a2}

\begin{lstlisting}[
	caption={ABBs present in both architectures (Kept)},
	label={lst:annex-kept-elements-query}
]
# Newly introduced prefix for this listing.
PREFIX ex: <http://www.example.org/eakg#>
# Case-study IDs: baseline = ex:id-47396; target = ex:id-46339.
SELECT ?shared_element
WHERE {
  ex:id-47396 archimate:Aggregation ?shared_element .
  ex:id-46339 archimate:Aggregation ?shared_element .
}
\end{lstlisting}

\begin{lstlisting}[
	caption={ABBs present only in the baseline (candidates for elimination)},
	label={lst:annex-eliminated-elements-query}
]
# Prefix used for the gap-analysis resources and properties.
PREFIX ex: <http://www.example.org/eakg#>
SELECT ?baseline_element
WHERE {
  ex:id-47396 archimate:Aggregation ?baseline_element .
  FILTER NOT EXISTS {
    ex:id-46339 archimate:Aggregation ?baseline_element .
  }
}
\end{lstlisting}

\begin{lstlisting}[
	caption={ABBs present only in the target (new gaps)},
	label={lst:annex-new-elements-query}
]
# Prefix used for the gap-analysis resources and properties.
PREFIX ex: <http://www.example.org/eakg#>
SELECT ?target_element
WHERE {
  ex:id-46339 archimate:Aggregation ?target_element .
  FILTER NOT EXISTS {
    ex:id-47396 archimate:Aggregation ?target_element .
  }
}
\end{lstlisting}

\begin{lstlisting}[
	caption={Classify elements as Kept, Eliminated, or New in one query},
	label={lst:annex-gap-classify-query}
]
# Newly introduced prefix for this query.
PREFIX ex: <http://www.example.org/eakg#>

SELECT ?ABB ?label ?state
WHERE {
    BIND(ex:id-46339 AS ?target)
    BIND(ex:id-47396 AS ?baseline)
    {
        ?baseline archimate:Aggregation ?ABB .
        ?target archimate:Aggregation ?ABB .
        ?ABB rdfs:label ?label .
        BIND("Kept" AS ?state)
    }
    UNION
    {
        ?baseline archimate:Aggregation ?ABB .
        FILTER NOT EXISTS {
            ?target archimate:Aggregation ?ABB .
        }
        ?ABB rdfs:label ?label .
        BIND("Eliminated" AS ?state)
    }
    UNION
    {
        ?target archimate:Aggregation ?ABB .
        FILTER NOT EXISTS {
            ?baseline archimate:Aggregation ?ABB .
        }
        ?ABB rdfs:label ?label .
        BIND("New" AS ?state)
    }
}
\end{lstlisting}

\begin{lstlisting}[
	caption={Mark gap decisions via SPARQL Update},
	label={lst:annex-gap-mark-decisions-query}
]
# Prefix used for the gap-analysis resources and properties.
PREFIX ex: <http://www.example.org/eakg#>

# Record one reviewed new element and one eliminated element.

INSERT {
  ?gap ex:gapNew ex:id-46919 .
  ?gap ex:gapEliminated ex:id-46864 .
}
WHERE {
  ?gap archimate:Association ex:id-47396 .
  ?gap archimate:Association ex:id-46339 .
}
\end{lstlisting}

\section{Architecture Evolution Query}
\label{sec:annex-archi-evo-queries}

\begin{lstlisting}[
	caption={SPARQL equivalent of a DL query for transformation lookup},
	label={lst:annex-equi-dl-tansformation}
]
# Prefix used for the transformation vocabulary.
PREFIX ex: <http://www.example.org/eakg#>

SELECT ?transformation
WHERE {
  ?transformation rdf:type ex:Transformation .
  ?project ex:name "Back-Office System Integration" .
  ?transformation ex:appliedBy ?project .
}
\end{lstlisting}

\section{Structural and Connectivity Smells Queries}
\label{sec:annex-structural-connectivity-smells-queries}

\begin{lstlisting}[
	caption={Detect cyclic dependencies},
	label={lst:annex-detect-cyclic-query}
]
SELECT DISTINCT ?node ?label
WHERE {
  ?node rdfs:label ?label .
  ?node archimate:DependencyRelationship+ ?node .
}
\end{lstlisting}

\begin{lstlisting}[
	caption={Detect shared persistency},
	label={lst:annex-detect-shared-persistency-query}
]
SELECT ?dataNode ?label
       (COUNT(DISTINCT ?businessNode) AS ?accessCount)
WHERE {
  ?dataNode a archimate:PassiveStructureElement .
  ?dataNode rdfs:label ?label .
  ?businessNode a archimate:BusinessLayer .
  ?businessNode archimate:Access ?dataNode .
}
GROUP BY ?dataNode ?label
HAVING (COUNT(DISTINCT ?businessNode) >= 2)
\end{lstlisting}

\begin{lstlisting}[
	caption={Detect data-only services},
	label={lst:annex-detect-data-only-query}
]
SELECT DISTINCT ?service ?label
WHERE {
  ?service a ?serviceType .
  FILTER (?serviceType IN (archimate:ApplicationService,
                           archimate:BusinessService))
  ?service rdfs:label ?label .
  ?relOut rdfs:subPropertyOf* archimate:Relationship .
  ?service ?relOut ?dataNode .
  ?dataNode a ?dataType .
  ?dataType rdfs:subClassOf* archimate:PassiveStructureElement .
  FILTER NOT EXISTS {
    ?relOther rdfs:subPropertyOf* archimate:Relationship .
    ?service ?relOther ?otherNode .
    ?otherNode a ?otherType .
    ?otherType rdfs:subClassOf* ?forbiddenClass .
    FILTER (?forbiddenClass IN (
      archimate:ActiveStructureElement,
      archimate:BehaviorElement))
  }
}
\end{lstlisting}

\section{Model Quality Smells Queries}
\label{sec:annex-model-quality-smells-queries}


\begin{lstlisting}[
	caption={Detect duplicate elements by name},
	label={lst:annex-detect-dup-element-query}
]
SELECT ?label (COUNT(DISTINCT ?node) AS ?instanceCount)
WHERE {
  ?node rdfs:label ?label .
  FILTER(isIRI(?node))
}
GROUP BY ?label
HAVING (COUNT(DISTINCT ?node) > 1)
\end{lstlisting}

\begin{lstlisting}[
	caption={Detect unnecessarily long documentation},
	label={lst:annex-detect-long-doc}
]
SELECT ?element ?comment
       (STRLEN(STR(?comment)) AS ?length)
WHERE {
  ?element rdfs:comment ?comment .
  FILTER (STRLEN(STR(?comment)) > 400)
}
\end{lstlisting}

\section{EA Smell Tagging Queries}

The following \texttt{INSERT} queries materialize \texttt{smells:hasSmell} triples on the elements or relations identified by the detection queries above. They use the additional prefix \texttt{smells: <http://purl.org/eakg/smells\#>}.

\begin{lstlisting}[
	caption={Tag cyclic dependency elements},
	label={lst:annex-tag-cyclic-dep-query}
]
INSERT {
  ?node smells:hasSmell smells:CyclicDependency .
} WHERE {
  ?node archimate:DependencyRelationship+ ?node .
}
\end{lstlisting}

\begin{lstlisting}[
	caption={Tag shared persistency data nodes},
	label={lst:annex-shared-persistency-query}
]
INSERT { ?dataNode smells:hasSmell smells:SharedPersistency . }
WHERE {
  {
    SELECT ?dataNode WHERE {
      ?dataNode a archimate:PassiveStructureElement .
      ?businessNode a archimate:BusinessLayer .
      ?businessNode archimate:Access ?dataNode .
    }
    GROUP BY ?dataNode
    HAVING (COUNT(DISTINCT ?businessNode) >= 2)
  }
}
\end{lstlisting}

\begin{lstlisting}[
	caption={Tag strict layers violation edges},
	label={lst:annex-tag-strict-layers-violation-query}
]
INSERT {
  ?relation smells:hasSmell smells:StrictLayersViolation .
} WHERE {
  {
    ?nodeA a archimate:BusinessLayer .
    ?nodeB a archimate:TechnologyLayer .
    ?nodeA ?relation ?nodeB .
  } UNION {
    ?nodeA a archimate:TechnologyLayer .
    ?nodeB a archimate:BusinessLayer .
    ?nodeA ?relation ?nodeB .
  }
  ?relation rdfs:subPropertyOf* archimate:Relationship .
  FILTER(STRSTARTS(STR(?relation), "http://www.example.org/eakg#id"))
}
\end{lstlisting}

\begin{lstlisting}[
	caption={Tag chatty service nodes},
	label={lst:annex-tag-chatty-service-nodes-query}
]
INSERT {
  ?service smells:hasSmell smells:ChattyService .
} WHERE {
  {
    SELECT ?service WHERE {
      ?service a ?serviceType .
      FILTER (?serviceType IN (archimate:ApplicationService,
                               archimate:BusinessService))
      { ?service ?relation ?connectedElement . }
      UNION
      { ?connectedElement ?relation ?service . }
      ?relation rdfs:subPropertyOf* archimate:Relationship .
      FILTER(STRSTARTS(STR(?relation),
             "http://www.example.org/eakg#id"))
    }
    GROUP BY ?service
    HAVING (COUNT(DISTINCT ?connectedElement) >= 10)
  }
}
\end{lstlisting}

\begin{lstlisting}[
	caption={Tag data-only service nodes},
	label={lst:annex-tag-data-only-query}
]
INSERT {
  ?service smells:hasSmell smells:DataService .
} WHERE {
  ?service a ?serviceType .
  FILTER (?serviceType IN (archimate:ApplicationService,
                           archimate:BusinessService))
  ?relOut rdfs:subPropertyOf* archimate:Relationship .
  ?service ?relOut ?dataNode .
  ?dataNode a ?dataType .
  ?dataType rdfs:subClassOf* archimate:PassiveStructureElement .
  FILTER NOT EXISTS {
    ?relOther rdfs:subPropertyOf* archimate:Relationship .
    ?service ?relOther ?otherNode .
    ?otherNode a ?otherType .
    ?otherType rdfs:subClassOf* ?forbiddenClass .
    FILTER (?forbiddenClass IN (
      archimate:ActiveStructureElement,
      archimate:BehaviorElement))
  }
}
\end{lstlisting}

\begin{lstlisting}[
	caption={Tag duplicate elements},
	label={lst:annex-tag-duplicate-query}
]
INSERT { ?node smells:hasSmell smells:Duplication . }
WHERE {
  ?node rdfs:label ?label .
  FILTER(isIRI(?node))
  {
    SELECT ?label WHERE {
      ?innerNode rdfs:label ?label .
      FILTER(isIRI(?innerNode))
    }
    GROUP BY ?label
    HAVING (COUNT(DISTINCT ?innerNode) > 1)
  }
}
\end{lstlisting}

\begin{lstlisting}[
	caption={Tag unnecessarily long documentation},
	label={lst:annex-tag-unecessary-long-doc-query}
]
INSERT { ?resource smells:hasSmell smells:UnnecessaryDocumentation . }
WHERE {
  ?resource rdfs:comment ?comment .
  FILTER (!STRSTARTS(STR(?resource), "http://purl.org/eakg/archimate"))
  FILTER (!STRSTARTS(STR(?resource), "http://purl.org/eakg/smells"))
  FILTER (STRLEN(STR(?comment)) > 400)
}
\end{lstlisting}

\section{Smell Tag Utility Queries}

Replace \texttt{smells:CyclicDependency} with any smell URI to parameterize these utility queries.

\begin{lstlisting}[
	caption={Reset (delete) tags for a specific smell},
	label={lst:annex-reset-smell}
]
DELETE WHERE {
  ?element smells:hasSmell smells:CyclicDependency .
}
\end{lstlisting}

\section{Visualization Queries}
\label{sec:annex-vissualization-queries}

\begin{lstlisting}[
	caption={Connection between Business and Application Layer},
	label={lst:annex-connection-business-application-query}
]
CONSTRUCT {
    ?layer1 ?rel1 ?layer2 .
    ?layer2 ?rel2 ?layer1 .
}
WHERE {
    # Select Classes
    ?layer1 rdfs:subClassOf archimate:ApplicationLayer .
    BIND(archimate:ApplicationLayer AS ?layer1Class)
    ?layer2 rdfs:subClassOf archimate:BusinessLayer .
    BIND(archimate:BusinessLayer AS ?layer2Class)
    ?a a ?layer1 .
    ?b a ?layer2 .
{
        ?a a ?layer1 .
        ?b a ?layer2 .
        ?a ?rel1 ?b .
    }
    UNION
    {
        ?a a ?layer1 .
        ?b a ?layer2 .
        ?b ?rel2 ?a .
    }
}
\end{lstlisting}

\begin{lstlisting}[
	caption={Dependency Chain Expansion Query},
	label={lst:dep-chain-expantion-query}
]
CONSTRUCT {
    ?node ?rel ?target .
}
WHERE {
    ?node ?rel ?target .
    ?rel rdfs:subPropertyOf ?parentRel .
    VALUES ?parentRel { 
        archimate:DependencyRelationship 
        archimate:StructuralRelationship 
    }
}
\end{lstlisting}

\begin{lstlisting}[
	caption={Relations between Technology Layer Elements},
	label={lst:relations-tech-layer-query}
]
CONSTRUCT {
    # 1. Draw the actual relationships (edges) between the nodes
    ?src ?rel ?target . 
    ?src rdfs:label ?srName .
    ?target rdfs:label ?targetName .
}
WHERE {
    ?src a archimate:TechnologyLayer ;
          ?rel ?target .
    ?target a archimate:TechnologyLayer .
}
\end{lstlisting}